\section{Background}
\subsection{Secure Encrypted Virtualization}

AMD \gls{sev} is a set of architectural extensions to AMD CPUs that enable the deployment of \glspl{cvm}.
\gls{sev} piggy-backs on the AMD \gls{sme}'s support for transparent memory encryption automatically performed by the \gls{sp}.
When \gls{sev} is enabled, each \gls{cvm} is provisioned a unique \gls{vek} that the \gls{sp} uses to encrypt its memory pages. A \gls{cvm}'s \gls{vek} is securely maintained by the CPU, and is indexed by the \gls{cvm}'s unique \gls{asid}.




\hdr{Address translation SEV}
Similar to traditional \glspl{vm}, \gls{cvm} in \gls{sev} also follows \gls{slat} provided by the hardware.
The MMU translates a virtual address accessed by software inside a \gls{cvm} in two levels. The first level is the translation between \gls{gva} and \gls{hpa}, which is . When a \gls{gpf}.

The page table also contains information about whether a page should be encrypted or not, in a bit in the PTE called the \texttt{c}-bit in both the \gls{gpt}'s and the \gls{npt}'s entries.
% To decide whether a page should be encrypted or not, \gls{sev} relies on the in .


% In the \gls{gpt}

\hdr{Security extensions}
Since \gls{sev}'s release, several vulnerabilities have been discovered by the research community.
For instance, the initial release of \gls{sev} does not protect the \gls{cvm} states that are stored in the \gls{vmcb}, allowing a malicious hypervisor to inject register values alter the \gls{cvm} execution or extract sensitive values that are stored in the registers~\cite{morbitzer_extracting_2019}.
Moreover, the lack of protection in the \gls{gpa} to \gls{hpa} mappings enables the hypervisor to replace the mapping or map a single \gls{gpa} to multiple \glspl{hpa}.
In response, AMD introduced two new security extensions.

\Gls{seves} introduce two key security features. First, all the registers stored in the virtual machine control block (VMCS) are automatically encrypted upon exiting a \gls{vm}, to prevent leakages. Second, Non-automatic Asynchronous Exit (NAE) and the \gls{ghcb} are introduced to allow the guest to selectively share states with the hypervisor when required.

The next security extension to \gls{sev} is \gls{sevsnp}, which adds integrity protection to encrypted pages.
It introduces a \gls{rmp} that tracks the metadata (e.g., owner of the page and the mapped \gls{gpa}) for each \gls{hpa}.
The processor then looks up the \gls{rmp} and checks for permission on every memory access.


\subsection{Side channels on CVMs}

\subsubsection{Nested Page Fault Side-channels}
% \begin{table}[t]
    \small
    \centering
    \begin{tabularx}{\columnwidth}{@{}lX@{}}
    \toprule
         \thead{Attack}& \thead{Leaked Pattern}  \\
         \midrule
         \cite{hetzelt2017security,werner2019severest}& System calls \\
         Se\cite{werner2019severest} & Application structure\\
         \cite{morbitzer2018severed,li2019exploiting,wilke2020sevurity,radev2020exploiting,li2021cipherleaks,li2022systematic} & Data accesses during critical window\\
         \cite{li2022systematic} & Function called during critical window \\
         \cite{li2022systematic}& Context switches\\
         \bottomrule
    \end{tabularx}
    \caption{Previously proposed attacks on SEV that requires the NPF controlled-channel}
    \label{tab:npf-attack}
\end{table}
% \begin{table}[t]
%     \small
%     \centering
%     \begin{tabularx}{\columnwidth}{@{}lccX@{}}
%     \toprule
%          \thead{Attack}& \multicolumn{2}{c}{\thead{Leaked Access Pattern}} &\thead{Implication} \\
%          & \thead{Code}  & \thead{Data}& \\ 
%          \midrule
%          \cite{hetzelt2017security,werner2019severest}&\\
%          \cite{werner2019severest} &\\
%          \cite{morbitzer2018severed,li2019exploiting,wilke2020sevurity,radev2020exploiting,li2021cipherleaks,li2022systematic} & \\ 
%          \cite{li2022systematic} & \\
%          \cite{li2022systematic}& \\
%          SEVerest~\cite{werner2019severest}& & &\\
%          Li et al.~\cite{li2019exploiting}&x&x& Build enc./dec. oracle \\
%          Li et al.~\cite{li2022systematic}& \\
%          \bottomrule
%     \end{tabularx}
%     \caption{Previously proposed attacks on SEV that requires the NPF controlled-channel}
%     \label{tab:npf-attack}
% \end{table}


The hypervisor can trigger \gls{npf} on any page mapped to a \gls{cvm} by unsetting the present bit (\texttt{p}-bit) in the \gls{npt}.
Later, when the victim \gls{cvm} accesses the targetted pages, a \texttt{\#NPF} is generated by the CPU, and the faulting address is written into the \gls{vmcb}.

Hence, by unsetting the \texttt{p}-bit of every page used by a \gls{cvm}, the hypervisor can monitor its memory accesses and code execution at the page level. By combining this coarse-grained information with prior knowledge about the target \gls{vm}, attackers can infer much more fine-grained information, such as the system call trace~\cite{hetzelt_security_2017}, or the sequence of functions that performs certain sensitive operations.


% For this reason, \thename focuses on mitigating the two most prominent side-channel attacks against \glspl{cvm}.


% There are two security impacts of the \gls{npf} controlled-channel.

% First, the hypervisor can peek into the \gls{cvm}'s execution states and potentially leak data.
% To this end, its impact is similar to the Page Fault controlled-channel on SGX~\cite{xu2015controlledchannel}, which has been shown to leak sensitive information.
% On SEV, researcher has been \gls{cvm}~\cite{werner2019severest}.
% TODO: We perform additional evaluations to confirm this.
% Second, leaking the execution states of the \gls{cvm} allows the attacker to employ more fine-grained attacks.
% We observe that most known attacks on \gls{sev} depend on the \gls{npf} side-channel.

% Li et al. \cite{li2019exploiting} uses the sequence of \gls{npf} to extract page accesses when serving SSH packet, then uses pattern matching to extract the page that contains the packet to be sent over SSH.

% \paragraph{Ciphertext Side-channels}
% The recent CipherLeak attack~\cite{li2021cipherleaks,li2022systematic} employs the \gls{npf} controlled-channel.


% \begin{table}[]
%     \centering
%     \begin{tabular}{ll}
%     \toprule
%          Side-channel & Leakage  \\
%          \midrule
%          Nested PF & Page-level read/write/execute access pattern  \\
%          Ciphertext & Cacheline-level write access pattern\\
%          \bottomrule
%     \end{tabular}
%     \caption{Caption}
%     \label{tab:my_label}
% \end{table}


\subsection{Unikernels}
Unikernels are specialized kernel images linked directly with the application to be executed on a single-address space environment. Due to them being highly specialized to the specific workload, applications running on unikernels achieve faster boot time, better I/O performance, and smaller image size~\cite{unikraft} compared to traditional kernels.


 \thename takes advantage of the absence of the user-kernel switch in unikernels to achieve efficient instrumentation on memory accesses.
We port an existing unikernel implementation, Unikraft~\cite{unikraft}, to work with AMD SEV-SNP to utilize its mature POSIX compatibility.


\subsection{Oblivious RAM (ORAM)}
\gls{oram} algorithms are well-studied and proven ways to eliminate discernable memory trace differences on a memory region.
Path ORAM~\cite{path-oram} is the most adapted algorithm due to its simplicity and efficiency.
Path \gls{oram} organizes data blocks into a tree structure. Accesses to data blocks within the trees are transformed into path access that seems random to an observer. 


% \thename utilizes an implementation of PathORAM~\cite{stefanov2018path} to manage the page pool.
